% GENERATED FILE. Edit canonical lessons, then run npm run generate:books.
% canonical-source-hash: fnv1a64:1618b5566dcba212
% canonical-lessons: PA-C339-inkar, PA-C339-talna-back-out, PA-C339-bahana, PA-C339-hargiz, PA-C339-arna

\chapter{No, Thank You}
\label{ch:pa-a2-no-thank-you}
\hlchaptermodality{339}

\begin{chapteropening}
\textbf{By the end of this chapter:} \emph{I can say no politely and give a reason.}

Complete the last lesson of chapter 339: I can say no politely and give a reason.
\end{chapteropening}

\section[\texorpdfstring{\pa{ਇਨਕਾਰ} (inkār)}{ਇਨਕਾਰ (inkār)}]{\pa{ਇਨਕਾਰ} (inkār) — a refusal; saying no}

\label{lesson:PA-C339-inkar}



\begin{quote}
Before the new one: say the Punjabi for an order, then the Punjabi for to interrupt.
\end{quote}

\subsection*{You'll want to know: \pa{ਇਨਕਾਰ}}
\textbf{\pa{ਇਨਕਾਰ}} — \emph{inkār} — \textquotedblleft{}a refusal; saying no\textquotedblright{}.

\begin{quote}
\pa{ਮੈਂ} \pa{ਇਨਕਾਰ} \pa{ਨਹੀਂ} \pa{ਕਰਦਾ} \pa{ਹਾਂ।}
\end{quote}

\begin{quote}
\emph{maiṁ inkār nahīṁ kardā hāṁ.}
\end{quote}

\begin{quote}
I am not refusing.
\end{quote}

A masculine noun.

\begin{cousinweb}
\textbf{\pa{ਇਨਕਾਰ}} (\emph{inkār}) is a loan, usually traced through Persian to Arabic \emph{inkār}, \textquotedblleft{}denial\textquotedblright{}.
\end{cousinweb}

\begin{grammarlens}[title={What you've built}]
One more word for asking and answering.
\end{grammarlens}

\subsection*{Your turn}
\begin{itemize}
\raggedright
  \item \emph{Say it:} \emph{inkār}
  \item \emph{Say it:} \emph{inkār}, once more
  \item \emph{Say it:} \emph{maiṁ inkār nahīṁ kardā hāṁ.} — I am not refusing.
  \item \emph{Recall:} say the Punjabi for an order, then the Punjabi for to interrupt, then say \emph{inkār} again
\end{itemize}

\begin{tcolorbox}[breakable,colback=teal!4,colframe=teal!35!black,title={Before you move on}]
What does \pa{ਇਨਕਾਰ} mean? (\textquotedblleft{}A refusal; saying no\textquotedblright{}.) Say it once more.
\end{tcolorbox}
\section[\texorpdfstring{\pa{ਟਲਣਾ} (ṭalṇā)}{ਟਲਣਾ (ṭalṇā)}]{\pa{ਟਲਣਾ} (ṭalṇā) — to back out; to go back on one's word}

\label{lesson:PA-C339-talna-back-out}



\begin{quote}
Before the new one: say the Punjabi for to interrupt, then the Punjabi for a refusal.
\end{quote}

\subsection*{You'll want to know: \pa{ਟਲਣਾ}}
\textbf{\pa{ਟਲਣਾ}} — \emph{ṭalṇā} — \textquotedblleft{}to back out; to go back on one's word\textquotedblright{}.

\begin{quote}
\pa{ਮੇਰਾ} \pa{ਦੋਸਤ} \pa{ਆਪਣੀ} \pa{ਗੱਲ} \pa{ਤੋਂ} \pa{ਨਹੀਂ} \pa{ਟਲਦਾ} \pa{ਹੈ।}
\end{quote}

\begin{quote}
\emph{merā dost āpṇī gall ton nahīṁ ṭaldā hai.}
\end{quote}

\begin{quote}
My friend does not go back on his word.
\end{quote}

Past (he …): \pa{ਉਹ} \pa{ਆਪਣੀ} \pa{ਗੱਲ} \pa{ਤੋਂ} \pa{ਟਲਿਆ}, \emph{oh āpṇī gall ton ṭaliā}.

\begin{cousinweb}
\textbf{\pa{ਟਲਣਾ}} (\emph{ṭalṇā}) is inherited from Sanskrit \emph{ṭalati}, \textquotedblleft{}moves aside\textquotedblright{} (a sense CDIAL infers); CDIAL lists Punjabi \emph{ṭalṇā}, \textquotedblleft{}to slip away\textquotedblright{}, under it (CDIAL 5450). CDIAL does not give the sense \textquotedblleft{}to back out\textquotedblright{}.
\end{cousinweb}

\begin{grammarlens}[title={What you've built}]
One more word for asking and answering.
\end{grammarlens}

\subsection*{Your turn}
\begin{itemize}
\raggedright
  \item \emph{Say it:} \emph{ṭalṇā}
  \item \emph{Say it:} \emph{ṭalṇā}, once more
  \item \emph{Say it:} \emph{merā dost āpṇī gall ton nahīṁ ṭaldā hai.} — My friend does not go back on his word.
  \item \emph{Recall:} say the Punjabi for to interrupt, then the Punjabi for a refusal, then say \emph{ṭalṇā} again
  \item \emph{Write it:} copy \textbf{\pa{ਟਲਣਾ}} once, checking the whole word
\end{itemize}

\begin{tcolorbox}[breakable,colback=teal!4,colframe=teal!35!black,title={Before you move on}]
What does \pa{ਟਲਣਾ} mean? (\textquotedblleft{}To back out; to go back on one's word\textquotedblright{}.) Say it once more.
\end{tcolorbox}
\section[\texorpdfstring{\pa{ਬਹਾਨਾ} (bahānā)}{ਬਹਾਨਾ (bahānā)}]{\pa{ਬਹਾਨਾ} (bahānā) — an excuse}

\label{lesson:PA-C339-bahana}



\begin{quote}
Before the new one: say the Punjabi for a refusal, then the Punjabi for to back out.
\end{quote}

\subsection*{You'll want to know: \pa{ਬਹਾਨਾ}}
\textbf{\pa{ਬਹਾਨਾ}} — \emph{bahānā} — \textquotedblleft{}an excuse\textquotedblright{}.

\begin{quote}
\pa{ਬਹਾਨੇ} \pa{ਨਾ} \pa{ਕਰੋ।}
\end{quote}

\begin{quote}
\emph{bahāne nā karo.}
\end{quote}

\begin{quote}
Don't make excuses.
\end{quote}

A masculine noun.

\begin{cousinweb}
\textbf{\pa{ਬਹਾਨਾ}} (\emph{bahānā}) is a loan, usually traced to Persian \emph{bahāna}, \textquotedblleft{}pretext\textquotedblright{}.
\end{cousinweb}

\begin{grammarlens}[title={What you've built}]
One more word for asking and answering.
\end{grammarlens}

\subsection*{Your turn}
\begin{itemize}
\raggedright
  \item \emph{Say it:} \emph{bahānā}
  \item \emph{Say it:} \emph{bahānā}, once more
  \item \emph{Say it:} \emph{bahāne nā karo.} — Don't make excuses.
  \item \emph{Recall:} say the Punjabi for a refusal, then the Punjabi for to back out, then say \emph{bahānā} again
  \item \emph{Write it:} copy \textbf{\pa{ਬਹਾਨਾ}} once, checking the whole word
\end{itemize}

\begin{tcolorbox}[breakable,colback=teal!4,colframe=teal!35!black,title={Before you move on}]
What does \pa{ਬਹਾਨਾ} mean? (\textquotedblleft{}An excuse\textquotedblright{}.) Say it once more.
\end{tcolorbox}
\section[\texorpdfstring{\pa{ਹਰਗਿਜ਼} (hargiz)}{ਹਰਗਿਜ਼ (hargiz)}]{\pa{ਹਰਗਿਜ਼} (hargiz) — ever; at all (with \pa{ਨਹੀਂ}: by no means)}

\label{lesson:PA-C339-hargiz}



\begin{quote}
Before the new one: say the Punjabi for to back out, then the Punjabi for an excuse.
\end{quote}

\subsection*{You'll want to know: \pa{ਹਰਗਿਜ਼}}
\textbf{\pa{ਹਰਗਿਜ਼}} — \emph{hargiz} — \textquotedblleft{}ever; at all (with \pa{ਨਹੀਂ}: by no means)\textquotedblright{}.

\begin{quote}
\pa{ਮੈਂ} \pa{ਉੱਥੇ} \pa{ਹਰਗਿਜ਼} \pa{ਨਹੀਂ} \pa{ਜਾਵਾਂਗਾ।}
\end{quote}

\begin{quote}
\emph{maiṁ utthe hargiz nahīṁ jāvāṅgā.}
\end{quote}

\begin{quote}
I will by no means go there.
\end{quote}

\begin{cousinweb}
\textbf{\pa{ਹਰਗਿਜ਼}} (\emph{hargiz}) is a loan, usually traced to Persian \emph{hargiz}, \textquotedblleft{}ever\textquotedblright{}.
\end{cousinweb}

\begin{grammarlens}[title={What you've built}]
One more word for asking and answering.
\end{grammarlens}

\subsection*{Your turn}
\begin{itemize}
\raggedright
  \item \emph{Say it:} \emph{hargiz}
  \item \emph{Say it:} \emph{hargiz}, once more
  \item \emph{Say it:} \emph{maiṁ utthe hargiz nahīṁ jāvāṅgā.} — I will by no means go there.
  \item \emph{Recall:} say the Punjabi for to back out, then the Punjabi for an excuse, then say \emph{hargiz} again
  \item \emph{Write it:} copy \textbf{\pa{ਹਰਗਿਜ਼}} once, checking the whole word
\end{itemize}

\begin{tcolorbox}[breakable,colback=teal!4,colframe=teal!35!black,title={Before you move on}]
What does \pa{ਹਰਗਿਜ਼} mean? (\textquotedblleft{}Ever; at all (with \pa{ਨਹੀਂ}: by no means)\textquotedblright{}.) Say it once more.
\end{tcolorbox}
\section[\texorpdfstring{\pa{ਅੜਨਾ} (aṛnā)}{ਅੜਨਾ (aṛnā)}]{\pa{ਅੜਨਾ} (aṛnā) — to dig one's heels in; to refuse to budge}

\label{lesson:PA-C339-arna}



\begin{quote}
Before the new one: say the Punjabi for an excuse, then the Punjabi for by no means.
\end{quote}

\subsection*{You'll want to know: \pa{ਅੜਨਾ}}
\textbf{\pa{ਅੜਨਾ}} — \emph{aṛnā} — \textquotedblleft{}to dig one's heels in; to refuse to budge\textquotedblright{}.

\begin{quote}
\pa{ਬੱਚਾ} \pa{ਅੜਦਾ} \pa{ਹੈ} \pa{ਅਤੇ} \pa{ਰੋਟੀ} \pa{ਨਹੀਂ} \pa{ਖਾਂਦਾ} \pa{ਹੈ।}
\end{quote}

\begin{quote}
\emph{bachchā aṛdā hai ate roṭī nahīṁ khāndā hai.}
\end{quote}

\begin{quote}
The child digs his heels in and won't eat his food.
\end{quote}

Past (he …): \pa{ਉਹ} \pa{ਆਪਣੀ} \pa{ਗੱਲ} \pa{ਉੱਤੇ} \pa{ਅੜਿਆ}, \emph{oh āpṇī gall utte aṛiā}.

\begin{cousinweb}
\textbf{\pa{ਅੜਨਾ}} (\emph{aṛnā}) is inherited from a reconstructed \emph{*aḍ}, \textquotedblleft{}obstruct, stop\textquotedblright{}; CDIAL lists Punjabi \emph{aṛnā}, \textquotedblleft{}to stop, stick fast\textquotedblright{}, under it (CDIAL 187).
\end{cousinweb}

\begin{grammarlens}[title={What you've built}]
One more word for asking and answering.
\end{grammarlens}

\subsection*{Your turn}
\begin{itemize}
\raggedright
  \item \emph{Say it:} \emph{aṛnā}
  \item \emph{Say it:} \emph{aṛnā}, once more
  \item \emph{Say it:} \emph{bachchā aṛdā hai ate roṭī nahīṁ khāndā hai.} — The child digs his heels in and won't eat his food.
  \item \emph{Recall:} say the Punjabi for an excuse, then the Punjabi for by no means, then say \emph{aṛnā} again
  \item \emph{Write it:} copy \textbf{\pa{ਅੜਨਾ}} once, checking the whole word
\end{itemize}

\begin{tcolorbox}[breakable,colback=teal!4,colframe=teal!35!black,title={Before you move on}]
What does \pa{ਅੜਨਾ} mean? (\textquotedblleft{}To dig one's heels in; to refuse to budge\textquotedblright{}.) Say it once more.
\end{tcolorbox}
